\chapter{Artin L 函数}\label{ch:04}

\begin{center}
Ehud de Shalit\footnote{原书脚注：“A disclaimer: This is an expository work based on a talk given at the instructional conference in Jerusalem in March 2001. None of the results---but all the mistakes---are due to the author.” 大意：本稿据 2001 年 3 月耶路撒冷教学会议上的报告写成，其中一切结果都不属于作者，而一切错误都属于作者。}
\end{center}

% 本章节号照原书（1--7），公式号 (n.m) 与之同源；章末恢复 book 类默认，免得影响后章。
\let\chFourOldTheSection\thesection
\renewcommand{\thesection}{\arabic{section}}

\section{引言}\label{c4:sec:1}

设 $\chi\colon(\mathbb{Z}/m\mathbb{Z})^{\times}\to\mathbb{C}^{\times}$ 为模 $m$ 的本原 Dirichlet 特征，并令 $K=\mathbb{Q}(\zeta)$，其中 $\zeta=e^{2\pi i/m}$。同构 $G=\operatorname{Gal}(K/\mathbb{Q})\cong(\mathbb{Z}/m\mathbb{Z})^{\times}$ 使我们能对 $\chi$ 附上一个特征 $\chi_{\mathrm{Gal}}\colon G\to\mathbb{C}^{\times}$，它满足
\begin{equation}\tag{1.1}\label{c4:eq:1.1}
\chi_{\mathrm{Gal}}(\sigma_p)=\chi(p),
\end{equation}
只要 $(p,m)=1$，而 $\sigma_p$ 是 $p$ 处的 Frobenius 自同构。所谓 $\sigma_p$，是 $p$ 在 $G$ 中分解群的标准生成元，它在 $K$ 的任一 $p$ 上素位的剩余域上诱导自同构 $x\mapsto x^{p}$。Kronecker--Weber 定理（Kronecker 1853，Weber 1886）断言：$G_{\mathbb{Q}}=\operatorname{Gal}(\overline{\mathbb{Q}}/\mathbb{Q})$ 的每个一维特征都形如 $\chi_{\mathrm{Gal}}$，其中 $\chi$ 适当选取。

于是 Dirichlet L 函数 $L(\chi,s)$ 可相应地写成
\begin{equation}\tag{1.2}\label{c4:eq:1.2}
L(\chi,s)=\prod_{p}\bigl(1-\chi_{\mathrm{Gal}}(\sigma_p)p^{-s}\bigr)^{-1}.
\end{equation}

\begin{annote}
$\sigma_p$ 只在 $p$ 上诸素位的分解群之间定到共轭，所以 (1.1) 的右端须与代表元的选取无关才良定；$\chi_{\mathrm{Gal}}$ 作为一维特征必是类函数，这一点自动成立。到非交换的 $\chi$，$L(\chi,s)$ 仍然良定，但理由变成 $\chi$ 是 $G$ 的类函数、$\sigma_v$ 的共轭类只依赖 $v$。(1.2) 的乘积实际上只对 $p\nmid m$ 写出，$p\mid m$ 处的因子留给第 3 节的惯性不变量补足；Kronecker--Weber 定理与 $\mathbb{Q}^{\mathrm{ab}}=\bigcup_{m}\mathbb{Q}(\zeta_m)$ 是同一件事的两种说法。
\par\vspace{0.15em}{\raggedleft 【来源：Serre《A Course in Arithmetic》Ch.~VI--VII（A）；Neukirch《Algebraic Number Theory》Ch.~VI（A）】\par}
\end{annote}

在 1923 年的一篇著名论文（“\"Uber eine neue Art von L-Reihen”）中，E.~Artin 把这一定义推广到 $G_{\mathbb{Q}}$ 的任意复值特征。

\section{有限群的表示}\label{c4:sec:2}

设 $G$ 为有限群。$G$ 的复值有限维表示构成的范畴 $\mathrm{Rep}_G$ 是半单的：每个表示都是不可约表示的直接和。若 $\rho\colon G\to GL(V)$ 是表示，其特征 $\chi=\chi_\rho$ 是函数
\begin{equation}\tag{2.1}\label{c4:eq:2.1}
\chi(g)=\operatorname{tr}\rho(g),\qquad g\in G .
\end{equation}
特征在同构意义下决定表示。这是正交关系的推论：不可约表示的特征 $\chi_1,\dots,\chi_s$ 构成 $G$ 上类函数空间的一组标准正交基，而类函数是在共轭类上取常值的函数，其内积为
\begin{equation}\tag{2.2}\label{c4:eq:2.2}
\langle u,v\rangle=|G|^{-1}\sum_{g\in G}u(g)\overline{v(g)} .
\end{equation}
因此不可约表示的（同构类）个数 $s$ 等于 $G$ 的共轭类个数。它们的维数 $n_i=\chi_i(1)$ 满足
\begin{equation}\tag{2.3}\label{c4:eq:2.3}
\sum_{i=1}^{s}n_i^{2}=|G| .
\end{equation}

\subsection{表示上的运算}\label{c4:sec:2.1}

\begin{itemize}[leftmargin=2.2em,itemsep=0.25em]
\item 直和：$\chi_{\rho_1\oplus\rho_2}=\chi_{\rho_1}+\chi_{\rho_2}$。
\item 张量积：$\chi_{\rho_1\otimes\rho_2}=\chi_{\rho_1}\chi_{\rho_2}$。
\item 反偶：$\chi_{\rho^{\vee}}=\overline{\chi_\rho}$，其中 $\rho^{\vee}(g)={}^{t}\rho(g^{-1})$。
\item 限制：$\chi_{\operatorname{Res}^{G}_{H}\rho}=\chi_\rho|_{H}$，其中 $\operatorname{Res}^{G}_{H}\rho\in\mathrm{Rep}_H$ 是 $\rho$ 到子群 $H$ 的限制。
\item 诱导：若 $\rho\colon H\to GL(W)$ 属于 $\mathrm{Rep}_H$，则 $R=\operatorname{Ind}^{G}_{H}\rho\colon G\to GL(V)$ 是一个表示，其空间为
\begin{equation}\tag{2.4}\label{c4:eq:2.4}
V=\{f\colon G\to W;\ f(hg)=\rho(h)f(g)\ \ \forall h\in H,\ g\in G\},
\end{equation}
而 $R$ 由右平移给出：
\begin{equation}\tag{2.5}\label{c4:eq:2.5}
(R(g)f)(g')=f(g'g).
\end{equation}
\end{itemize}

$\chi_R=\operatorname{Ind}^{G}_{H}\chi_\rho$ 的公式略繁：
\begin{equation}\tag{2.6}\label{c4:eq:2.6}
(\operatorname{Ind}^{G}_{H}\chi)(g)=\sum_{\gamma\in G/H,\ \gamma^{-1}g\gamma\in H}\chi(\gamma^{-1}g\gamma).
\end{equation}

诱导与限制关于 $G$ 与 $H$ 上的内积互为对偶，这就是 Frobenius 互反律：
\begin{equation}\tag{2.7}\label{c4:eq:2.7}
\langle\psi,\chi|_{H}\rangle_{H}=\langle\operatorname{Ind}^{G}_{H}\psi,\chi\rangle_{G} .
\end{equation}

Frobenius 互反律的另一个推论（也不难直接证明）是“投影公式”
\begin{equation}\tag{2.8}\label{c4:eq:2.8}
\operatorname{Ind}^{G}_{H}\bigl(\psi\cdot\operatorname{Res}^{G}_{H}\chi\bigr)=\operatorname{Ind}^{G}_{H}(\psi)\cdot\chi .
\end{equation}

\begin{annote}
(2.4) 的函数方程写在左边，即 $f(hg)=\rho(h)f(g)$，而 (2.5) 的表示却由右平移给出；这套左右约定与 $G/H$ 的陪集代表元选取配套，抄写时最容易出错的是把 $f(gh)$ 与 $f(hg)$ 混用。由此，(2.6) 在 $g$ 不共轭进入 $H$ 时求和为空、特征取 $0$，这正是诱导特征与限制特征的对偶性 (2.7) 能够成立的原因。另一处常见疏漏是忘记 (2.6) 中的 $\gamma$ 取遍 $G/H$ 的代表元，而不是取遍整个 $G$。
\par\vspace{0.15em}{\raggedleft 【来源：Serre《Linear Representations of Finite Groups》§7.1--7.2（A）；Isaacs《Character Theory of Finite Groups》Ch.~5（A）】\par}
\end{annote}

（右）正则表示是 $\operatorname{Ind}^{G}_{1}1$，即 $G$ 在 $G$ 上复值函数空间里由右平移给出的表示。由 Frobenius 互反律，正则表示含每个不可约表示，重数等于该表示的维数。形如 $\operatorname{Ind}^{G}_{H}\rho$ 的表示（其中 $\rho$ 一维）称为单项的。

虚拟表示 $\rho$ 是不可约表示的整系数形式和。其特征可按显然的方式定义，并且仍然决定 $\rho$。虚拟表示（等价地，其特征）关于直和与张量积——等价地，类函数的通常加法与乘法——构成交换环 $R_G$。限制（相应地，诱导）函子延拓为环（相应地，群）同态 $R_G\to R_H$（相应地，$R_H\to R_G$），Frobenius 互反律与投影公式仍成立。

\begin{annote}
$R_G$ 不只是加法群而是环：乘法由张量积给出，(2.8) 说明 $\operatorname{Ind}$ 的像是理想。Brauer 定理（定理 4.2）断言 $R_G$ 由单项特征生成，但生成系数只能取整数。第 4 节会看到，这个“只能取整数”正是 Brauer 方法只给出亚纯延拓的原因；系数若能全取非负，Artin 猜想立刻成立。
\par\vspace{0.15em}{\raggedleft 【来源：Serre 前引书 Ch.~10（A）；Isaacs 前引书 Ch.~8（A）】\par}
\end{annote}

{\small
\par\medskip\noindent\textbf{参考文献}\par\nopagebreak
\noindent[1] J.-P. Serre, \emph{Linear Representations of Finite Groups}, Graduate Texts in Mathematics 42, Springer-Verlag, New York, 1977.\par
}

\section{Artin L 函数}\label{c4:sec:3}

\subsection{整体域的 Galois 群}\label{c4:sec:3.1}

设 $k$ 为整体域。于是 $k$ 或是 $\mathbb{Q}$ 的有限扩张（数域），或是有限域 $\mathbb{F}_q$ 上光滑、射影、几何不可约曲线 $X$ 的函数域（函数域）。设 $K$ 是 $k$ 的有限 Galois 扩张，$G=\operatorname{Gal}(K/k)$。

若 $v$ 是 $k$ 的一个素位（数域情形为非零素理想，函数域情形为 $X$ 上的素除子），$\bar v$ 是 $K$ 上位于 $v$ 之上的素位，则记 $G_{\bar v}$ 为 $\bar v$ 在 $G$ 中的分解群，$I_{\bar v}$ 为惯性群。记 $\sigma_{\bar v}$ 为 $G_{\bar v}/I_{\bar v}$ 的 Frobenius 自同构，并记 $\sigma_v$ 为诸 $\sigma_{\bar v}$（$\bar v\mid v$）中的任一个，它们彼此共轭。

数域情形，设 $v$ 是 $k$ 的阿基米德位，$\bar v$ 是 $K$ 上位于其上的位。若二者皆实或皆复，令 $G_{\bar v}=\{1\}$。若 $v$ 实而 $\bar v$ 复，令 $G_{\bar v}=I_{\bar v}$ 为 2 阶群，它由唯一的 $\sigma_{\bar v}\in G$ 生成，而 $\sigma_{\bar v}$ 交换对应于 $\bar v$ 的两个（共轭的）复嵌入。

使 $I_{\bar v}\neq1$ 的素位（或阿基米德位）$v$ 称为分歧的。分歧的素位只有有限多个。

\subsection{Artin L 函数}\label{c4:sec:3.2}

设 $\rho\colon G\to GL(V)$ 为表示，$\chi$ 为其特征。对每个素位 $v$，商 $G_{\bar v}/I_{\bar v}$ 作用在 $V^{I_{\bar v}}$ 上，故下式良定：
\begin{equation}\tag{3.1}\label{c4:eq:3.1}
L(\chi,s)=\prod_{v}L_v(\chi,s),
\end{equation}
其中乘积遍历所有素位，$v$ 处的 Euler 因子为
\begin{equation}\tag{3.2}\label{c4:eq:3.2}
L_v(\chi,s)=\det\bigl(1-\sigma_vq_v^{-s}\mid V^{I_{\bar v}}\bigr)^{-1}
\end{equation}
（$q_v$ 是 $v$ 处剩余域的基数）。Euler 因子作为 $q_v^{-s}$ 的多项式，其次数为 $\dim V^{I_{\bar v}}$；对几乎所有 $v$，它等于 $\dim\rho$。若 $V^{I_{\bar v}}\neq V$，则称素位 $v$ 在 $\rho$ 中分歧（或“坏”）。若 $\rho$ 忠实，这与 $v$ 在 $K/k$ 中分歧是一回事。

坏素位处的 Euler 因子，除了自然之外，是由“函数方程好看”这一愿望以及诸如关于 $\chi$ 的可加性（见下）这类性质所决定的。若省略它们，它们会以补丁因子的形式在函数方程中重新出现。若 $\rho$ 一维，坏素位处 $V^{I_{\bar v}}=0$，故 Euler 积只覆盖好素位；但一般情形坏 Euler 因子可能不等于 $1$。

\begin{annote}
坏素位处的因子由惯性不变量 $V^{I_{\bar v}}$ 定义，这是使 Euler 积在 $\operatorname{Re}(s)>1$ 收敛、又使函数方程整齐的唯一自然选择。若在坏素位处只写 $1$，函数方程右端就多出只依赖分歧数据的因子，书中称之为补丁因子。一维情形 $V^{I_{\bar v}}=0$ 使坏因子等于 $1$，所以 Dirichlet L 函数看不出这一层；一般情形 $\dim V^{I_{\bar v}}$ 在分歧处严格小于 $\dim\rho$，差值正是导子要记录的信息。
\par\vspace{0.15em}{\raggedleft 【来源：Tate《Local constants》，Durham 会议录（A）；Serre《Zeta and L functions》，Schilling 编（A）】\par}
\end{annote}

等式
\begin{equation}\tag{3.3}\label{c4:eq:3.3}
\log\det(1-A)^{-1}=\sum_{n=1}^{\infty}\frac{\operatorname{tr}(A^{n})}{n}
\end{equation}
（$A$ 为矩阵）给出下面这个有用的公式。令
\begin{equation}\tag{3.4}\label{c4:eq:3.4}
\chi(\sigma_v^{n})=|I_{\bar v}|^{-1}\sum_{\sigma\bmod I_{\bar v}=\sigma_v^{n}}\chi(\sigma)
\end{equation}
为 $\chi$ 在陪集 $\sigma_v^{n}$ 上的平均值。则
\begin{equation}\tag{3.5}\label{c4:eq:3.5}
\log L_v(\chi,s)=\sum_{n=1}^{\infty}\frac{\chi(\sigma_v^{n})}{nq_v^{ns}} .
\end{equation}

Euler 积在 $\operatorname{Re}(s)>1$ 时收敛，且在紧集上一致收敛，故它在该区域表示一个无处为零的解析函数，并且在每个形如 $\operatorname{Re}(s)\ge1+\varepsilon$ 的闭半平面内有界。

\subsection{函数域情形}\label{c4:sec:3.3}

（参看所引 Serre 的综述 [3]。）此情形下 Artin L 函数的解析性质已由 Weil 关于有限域上曲线的工作弄清楚，并对数域情形下应当期待什么提供了启示。设 $\mathbb{F}_q$ 在 $K$ 中代数闭，即扩张 $K/k$ 不涉及常数域扩张。此时 $K$ 是 $\mathbb{F}_q$ 上唯一的光滑射影几何不可约曲线 $Y$ 的函数域，而 $K/k$ 对应一个以 $G$ 为覆盖群的有限 Galois 覆盖 $Y\to X$。记 $\bar Y$ 与 $\bar X$ 为 $\overline{\mathbb{F}}_q$ 上相应的曲线。设 $G$ 的表示 $\rho$ 不可约且非平凡，否则我们研究的只是 $k$ 的 zeta 函数 $\zeta_k(s)$。

回忆 Jacobian 簇 $\operatorname{Jac}(\bar Y)$ 是群簇，其 $\overline{\mathbb{F}}_q$-有理点是 $\bar Y$ 上次数 $0$ 的除子类模主除子。若 $\ell$ 是不同于 $p=\operatorname{char}(\mathbb{F}_q)$ 的素数，其 $\ell^{r}$-挠点为
\begin{equation}\tag{3.6}\label{c4:eq:3.6}
\operatorname{Jac}(\bar Y)[\ell^{r}]\cong(\mathbb{Z}/\ell^{r}\mathbb{Z})^{2g},
\end{equation}
其中 $g=\operatorname{genus}(\bar Y)$。乘以 $\ell$ 之下这些群构成逆向系统，其投射极限称为 $\operatorname{Jac}(\bar Y)$ 的 $\ell$ 进 Tate 模。记 $H^{1}(\bar Y,\mathbb{Q}_\ell)$ 为其 $\mathbb{Q}_\ell$-对偶，
\begin{equation}\tag{3.7}\label{c4:eq:3.7}
H^{1}(\bar Y,\mathbb{Q}_\ell)=\Hom\bigl(\varprojlim\operatorname{Jac}(\bar Y)[\ell^{r}],\mathbb{Q}_\ell\bigr).
\end{equation}
它是 $\mathbb{Q}_\ell$ 上 $2g$ 维向量空间，带有 $G$ 与 $\operatorname{Gal}(\overline{\mathbb{F}}_q/\mathbb{F}_q)$ 的交换作用。记 $H^{1}(\bar Y,\mathbb{Q}_\ell)(\rho)$ 为 $\rho$-等型分量。为定义它，先把 $\rho$ 看作 $\mathbb{Q}_\ell$ 上向量空间 $V$ 中的表示，而非复数上的表示，然后定义
\begin{equation}\tag{3.8}\label{c4:eq:3.8}
H^{1}(\bar Y,\mathbb{Q}_\ell)(\rho)=\Hom_G\bigl(V,H^{1}(\bar Y,\mathbb{Q}_\ell)\otimes\overline{\mathbb{Q}}_\ell\bigr).
\end{equation}
由于 $G$ 与 $\operatorname{Gal}(\overline{\mathbb{F}}_q/\mathbb{F}_q)$ 的作用交换，此空间仍带有 $\operatorname{Gal}(\overline{\mathbb{F}}_q/\mathbb{F}_q)$ 的作用。

设 $F$ 为 $q$ 次幂的 Frobenius 自同构，它是 $\operatorname{Gal}(\overline{\mathbb{F}}_q/\mathbb{F}_q)$ 的拓扑生成元。上一节记录的 $\log L_v(\chi,s)$ 的表达式，使我们能把 $L(\chi,s)$ 用 $gF^{n}$ 在 $Y(\overline{\mathbb{F}}_q)$ 上的不动点个数 $\Lambda(gF^{n})$ 表出，其中 $n=1,2,\dots$、$g\in G$：
\[
\log L(\chi,s)=\sum_{n=1}^{\infty}\Bigl\{|G|^{-1}\sum_{g\in G}\chi(g^{-1})\Lambda(gF^{n})\Bigr\}\frac{q_v^{-ns}}{n}.
\]
这是初等的，只涉及一些代数。注意 $gF^{n}$ 的每个不动点必落在 $X(\mathbb{F}_{q^{n}})$ 的某个点之上，因为 $G$ 作用在 $Y\to X$ 的纤维上。

A.~Weil 的理论（\emph{Courbes algébriques et variétés abéliennes}，1948），特别是其中 Lefschetz 不动点公式的类比，现在蕴涵
\begin{equation}\tag{3.9}\label{c4:eq:3.9}
L(\chi,s)=Z(X,\rho,q^{-s}),
\end{equation}
其中
\begin{equation}\tag{3.10}\label{c4:eq:3.10}
Z(X,\rho,t)=\det\bigl(1-Ft\mid H^{1}(\bar Y,\mathbb{Q}_\ell)(\rho)\bigr)
\end{equation}
是 $\mathbb{Q}$ 上 $t$ 的多项式。

于是得到下列性质中的第一条，即解析延拓。第二条即函数方程，来自 Poincaré 对偶，也就是 Weil 配对的存在性及其关于 $G$ 的函子性。第三条是函数域的 Riemann 假设，为 Weil 的主要成就之一，要深刻得多。但请注意：由 (i) 以及把 Artin L 级数之积与 zeta 函数联系起来的公式（见下），$L(\chi,s)$ 的每个零点都作为 $\zeta_K(s)$ 的零点出现，故 Artin L 函数的 RH 由 zeta 函数的 RH 推出。

\par\medskip\noindent\textbf{定理 3.1（A.~Weil）}\quad 设 $\rho$ 不可约且非平凡。则有：
\begin{enumerate}[label=(\roman*),leftmargin=2.6em,itemsep=0.3em]
\item Artin L 级数 $L(\chi,s)$ 是 $q^{-s}$ 的多项式，因而在整个复平面上解析。
\item 有函数方程
\begin{equation}\tag{3.11}\label{c4:eq:3.11}
Z(X,\rho,t)=\epsilon Z(X,\rho^{\vee},q^{-1}t^{-1}),
\end{equation}
等价地说，它把 $L(\chi,1-s)$ 与 $L(\tilde\chi,s)$ 联系起来；其中
\[
\epsilon=at^{d(\rho)},\qquad d(\rho)=\dim H^{1}(\bar Y,\mathbb{Q}_\ell)(\rho),\qquad
a=(-1)^{d(\rho)}\det\bigl(F\mid H^{1}(\bar Y,\mathbb{Q}_\ell)(\rho)\bigr).
\]
\item $L(\chi,s)$ 的所有零点都在直线 $\operatorname{Re}(s)=1/2$ 上。
\end{enumerate}

\begin{annote}
(3.11) 中的 $\epsilon$ 不是常数，而是含 $t^{d(\rho)}$ 的单项式；把 $t^{d}$ 并入根数写成 $at^{d(\rho)}$，是 Weil 的写法。由 Weil 的结果，$F$ 在 $H^{1}$ 上的特征值模长为 $q^{1/2}$，故 $|a|=q^{d(\rho)/2}$；于是在中心线 $t=q^{-1/2}$ 处 $|\epsilon|=1$，函数方程在 $s\leftrightarrow1-s$ 下保持模长。$a$ 中的符号 $(-1)^{d(\rho)}$ 来自 $\det(1-Ft)$ 的常数项。
\par\vspace{0.15em}{\raggedleft 【来源：Weil《Courbes algébriques et variétés abéliennes》（A）；Milne《Étale Cohomology》Ch.~VI（A）】\par}
\end{annote}

这些结果被 Dwork、Grothendieck 与 Deligne 用（Grothendieck 与 Deligne 的）$\ell$ 进平展上同调推广到任意维有限域上的簇。数域情形中相应的问题属于数论最引人入胜的猜想之列。

\subsection{$L(\chi,s)$ 的初步性质}\label{c4:sec:3.4}

\begin{itemize}[leftmargin=2.2em,itemsep=0.3em]
\item 可加性：$L(\chi_1+\chi_2,s)=L(\chi_1,s)L(\chi_2,s)$。
\item 诱导：若 $K/k'/k$ 是中间扩张，$\chi$ 是 $H=\operatorname{Gal}(K/k')$ 的某个表示的特征，则
\begin{equation}\tag{3.12}\label{c4:eq:3.12}
L\bigl(\operatorname{Ind}^{G}_{H}\chi,s\bigr)=L(\chi,s).
\end{equation}
注意左端的底域是 $k$，右端是 $k'$。
\item 膨胀：若 $K'/K/k$ 是更大的扩张，且 $K'$ 在 $k$ 上也是 Galois 的，$\chi'$ 是 $G'=\operatorname{Gal}(K'/k)$ 的特征，它由先投影到 $G$、再对 $\chi$ 取值而得，则
\begin{equation}\tag{3.13}\label{c4:eq:3.13}
L(\chi',s)=L(\chi,s).
\end{equation}
\item 类域论：设 $G$ 交换、$\chi$ 不可约，即一维。类域论对 $\chi$ 配一个广义理想类群 $\mathrm{Cl}(\mathfrak{f})$ 上的特征 $\chi_{\mathrm{CFT}}$。数域情形，$\mathrm{Cl}(\mathfrak{f})$ 由与 $\mathfrak{f}$ 互素的理想，模掉由全正且 $\equiv1\bmod\mathfrak{f}$ 的元素生成的主理想而得；函数域情形则是用 $X$ 的某个广义 Jacobian 构造的类似群。例如，若 $K$ 是 $k$ 的 Hilbert 类域，即 $K/k$ 无分歧，则可取 $\mathfrak{f}=1$，$\mathrm{Cl}(\mathfrak{f})$ 就是 $k$ 的（窄）类群。这个对应由
\begin{equation}\tag{3.14}\label{c4:eq:3.14}
\chi_{\mathrm{CFT}}([v])=\chi(\sigma_v)
\end{equation}
确定。
\end{itemize}

\begin{annote}
诱导公式 (3.12) 两端的底域不同：左端是 $k$ 上的 Artin L 函数，右端是中间域 $k'$ 上的。把 $k'$ 误当 $k$，导子与判别式都会算错，函数方程里的 $|d_k|$ 也会用错。膨胀 (3.13) 走的是反方向，特征先经 $G'\to G$ 拉回再取值。两条合起来说明 $L(\chi,s)$ 只依赖 $\chi$ 与底域，与 Galois 群怎样实现无关。
\par\vspace{0.15em}{\raggedleft 【来源：Martinet《Character theory and Artin L-functions》，Durham 会议录（A）；Serre《Zeta and L functions》（A）】\par}
\end{annote}

两件事——(i) 如此定义的 $\chi_{\mathrm{CFT}}$ 过 $\mathrm{Cl}(\mathfrak{f})$ 分解；(ii) $\chi\mapsto\chi_{\mathrm{CFT}}$ 是双射——就是类域论的一言之总结。类域论在 20 世纪上半叶发展起来，有许多好书以不同途径讲述它，解析的与上同调的、局部的与整体的皆有。我们在此不深入这一主题。

由定义，
\begin{equation}\tag{3.15}\label{c4:eq:3.15}
L(\chi,s)=L(\chi_{\mathrm{CFT}},s),
\end{equation}
而后者当 $\chi$ 非平凡时有解析延拓与函数方程（Hecke）。

引言中的对应 $\chi_{\mathrm{Gal}}\leftrightarrow\chi$ 是此处讨论的对应 $\chi\leftrightarrow\chi_{\mathrm{CFT}}$ 的特例。一般地，类域论就是“$\mathrm{GL}(1)$ 的自守表示论”；而 Galois 特征与理想类群特征之间的对应，是 Langlands 关于 motivic L 函数与自守 L 函数的猜想的第一例，也是其动机。

\par\medskip\noindent\textbf{例}\quad 前两条性质与正则表示的著名分解的一个推论是
\begin{equation}\tag{3.16}\label{c4:eq:3.16}
\zeta_K(s)=\zeta_k(s)\prod_{\chi\neq1}L(\chi,s)^{\chi(1)} .
\end{equation}

今后我们专注于数域情形。

{\small
\par\medskip\noindent\textbf{参考文献}\par\nopagebreak
\noindent[1] J. Martinet, Character theory and Artin L-functions, in \emph{Algebraic Number Fields: Proceedings of the Durham Symposium}, Academic Press, New York, 1977.\par
\noindent[2] J. Tate, \emph{Les conjectures de Stark sur les fonctions L d'Artin en $s=0$: Notes d'un cours a Orsay}, Progress in Mathematics 47, Birkh\"auser, Basel, 1984, Chapter 0.\par
\noindent[3] J.-P. Serre, Zeta and L functions, in O.~F.~G. Schilling, ed., \emph{Arithmetic Algebraic Geometry}, Harper and Row, New York, 1965.\par
}

\section{Brauer 定理、亚纯延拓与函数方程}\label{c4:sec:4}

\subsection{阿基米德 Euler 因子}\label{c4:sec:4.1}

众所周知，而且最先在乘积公式中见到：在数域与函数域的类比里，数域的有限素位对应开曲线，而为得到完备曲线的类比还须把阿基米德位一并考虑。对 L 函数的后果（Riemann zeta 函数已开其端）是：数域的 L 函数应在阿基米德位处补上 Euler 因子，才能得到好看的函数方程。这些因子的确切形状是从大量特例的研究中得来的。只有 Tate 的论文（将由 S.~Kudla 讨论）从第一原理出发，在交换 L 函数的特殊情形下对有限与无限 Euler 因子给出了统一推导。如今，Hodge 理论的一般原理使人们能预测这些 Euler 因子的正确形状，即使解析延拓与函数方程仍然未知。下面给出 Artin L 函数的情形。

定义
\begin{equation}\tag{4.1}\label{c4:eq:4.1}
\Gamma_{\mathbb{R}}(s)=\pi^{-s/2}\Gamma(s/2),\qquad
\Gamma_{\mathbb{C}}(s)=\Gamma_{\mathbb{R}}(s)\Gamma_{\mathbb{R}}(s+1)=2(2\pi)^{-s}\Gamma(s),
\end{equation}
并记 $\Gamma_v(s)$ 为 $\Gamma_{\mathbb{R}}(s)$ 或 $\Gamma_{\mathbb{C}}(s)$，视 $v$ 为实或复而定。令
\begin{equation}\tag{4.2}\label{c4:eq:4.2}
n_+(\chi,v)=\dim V^{G_{\bar v}},\qquad n_-(\chi,v)=\dim V/V^{G_{\bar v}},
\end{equation}
以及
\begin{equation}\tag{4.3}\label{c4:eq:4.3}
L_v(\chi,s)=\Gamma_v(s)^{n_+(\chi,v)}\Gamma_v(s+1)^{n_-(\chi,v)}.
\end{equation}

\begin{annote}
$\Gamma_{\mathbb{R}}$ 与 $\Gamma_{\mathbb{C}}$ 是 Tate 论文中局部 zeta 积分在实位与复位处的显式值。$\Gamma_{\mathbb{R}}(s)=\pi^{-s/2}\Gamma(s/2)$ 对应 $\mathbb{R}^{\times}$ 上的乘法特征 $|x|^{s}$，$\Gamma_{\mathbb{C}}$ 则是两个 $\Gamma_{\mathbb{R}}$ 之积。(4.3) 用 $n_+=\dim V^{G_{\bar v}}$ 与 $n_-=\dim V/V^{G_{\bar v}}$ 把 $s$ 与 $s+1$ 分配下去：复位的 $G_{\bar v}$ 平凡，$n_+=n_-=\dim V$，于是 $\Gamma_{\mathbb{C}}(s)^{\dim V}=2^{\dim V}(2\pi)^{-s\dim V}\Gamma(s)^{\dim V}$，与经典的 $2(2\pi)^{-s}\Gamma(s)$ 因子逐项吻合。
\par\vspace{0.15em}{\raggedleft 【来源：Tate《Fourier analysis in number fields and Hecke's zeta-functions》（A）；Bump《Automorphic Forms and Representations》Ch.~3（A）】\par}
\end{annote}

\subsection{Artin 导子}\label{c4:sec:4.2}

Artin 导子 $\mathfrak{f}(\chi)$ 是 $k$ 的一个理想，它度量 $\rho$ 的分歧。例如，若 $\chi$ 是模 $m$ 的本原 Dirichlet 特征，其导子就是 $m$。Artin 导子恰好被在 $\rho$ 中分歧的素位整除。对 $k$ 的任一素位 $v$，令 $I_{\bar v}=G_0\supset G_1\supset G_2\supset\cdots$ 为 $\bar v/v$ 的高阶分歧群序列。令
\begin{equation}\tag{4.4}\label{c4:eq:4.4}
f(\chi,v)=|G_0|^{-1}\sum_{i=0}^{\infty}|G_i|\operatorname{codim}V^{G_i}.
\end{equation}
（和为有限，且对好的 $v$ 为零。）下面是一个深刻的结果，属于 Artin：上式是整数；在交换情形它归结为局部类域论中的 Hasse--Arf 定理。令
\begin{equation}\tag{4.5}\label{c4:eq:4.5}
\mathfrak{f}(\chi)=\prod_{v}v^{f(\chi,v)} .
\end{equation}

\begin{annote}
(4.4) 取整数值是 Artin 的深刻定理，它等价于分歧群的跳跃点在上编号下受整性约束；交换情形即局部类域论的 Hasse--Arf 定理。$|G_i|$ 随 $i$ 递减而 $\operatorname{codim}V^{G_i}$ 递增，两者相抵才使级数有限。对好的 $v$ 有 $G_0=1$、求和为零，故 $\mathfrak{f}(\chi)$ 恰被分歧素位整除，这与“导子度量分歧”的说法一致。
\par\vspace{0.15em}{\raggedleft 【来源：Serre《Local Fields》Ch.~VI（A）；Tate《Local constants》，Durham 会议录（A）】\par}
\end{annote}

\subsection{亚纯延拓与函数方程}\label{c4:sec:4.3}

设 $d_k$ 为 $k$ 的判别式，定义
\begin{equation}\tag{4.6}\label{c4:eq:4.6}
\Lambda(\chi,s)=\bigl(|d_k|^{\chi(1)}N\mathfrak{f}(\chi)\bigr)^{s/2}\prod_{v\mid\infty}L_v(\chi,s)\cdot L(\chi,s).
\end{equation}
这是增广的 L 级数。可以验证它关于 $\chi$ 满足与 $L(\chi,s)$ 相同的函子性，即可加性、诱导与膨胀。这需要 $\Gamma$ 函数的简单公式，还需要知道 Artin 导子与判别式在域塔中的行为。

\par\medskip\noindent\textbf{定理 4.1}\quad $\Lambda(\chi,s)$ 亚纯延拓到整个复平面。它满足函数方程
\begin{equation}\tag{4.7}\label{c4:eq:4.7}
\Lambda(\chi,1-s)=W(\chi)\Lambda(\tilde\chi,s),
\end{equation}
其中 $W(\chi)\in\mathbb{C}^{\times}$ 是绝对值 $1$ 的常数，称为 Artin 根数。

Artin 导子在函数域情形也有定义，且可以证明上面的公式与函数域情形函数方程的精确形式相类似。然而与那里相反，此处没有直接的途径，也没有（无穷维的？）“算术上同调”能像 Weil 的理论在函数域那样解释这个定理，尽管寻找这样的理论是一个很热的课题。定理的证明分两步。

\textbf{第一步。}用类域论（“$\mathrm{GL}(1)$ 的自守表示论”）证明 $\rho$ 一维的情形。这属于经典，归功于 Dirichlet（$\mathbb{Q}$ 上）、Hecke（一般情形），以及 Iwasawa 与 Tate（用伊代尔群上调和分析的现代途径）。我们不再深入，因为这是 Kudla 讲 Tate 论文的那一章的主题。只注明一点：与特征 $\chi_{\mathrm{CFT}}$ 打交道，是在仅由 $k$ 构造的对象上打交道，而非在先验未知的 Galois 群上打交道；这使证明归结为若干具体可及的对象——$k$ 嵌入欧氏空间时的几何考虑（Hecke 的途径），或由 $k$ 及其完备化构造的局部紧群上的调和分析（Tate 的途径）。

\textbf{第二步。}用有限群表示论中 Brauer 的一个定理，把一般情形归到一维 $\chi$ 的情形。

\subsection{Brauer 诱导定理}\label{c4:sec:4.4}

\par\medskip\noindent\textbf{定理 4.2（R.~Brauer，1947）}\quad 设 $G$ 为有限群。则 $G$ 的每个特征都是单项特征的 $\mathbb{Z}$-线性组合。

回忆单项特征形如 $\operatorname{Ind}^{G}_{H}\chi$，其中 $\chi$ 一维。一般不能保证所有系数为正，即使允许有理系数；而正如我们将看到的，这正是 Artin 猜想困难的根源。确实，若 $\chi$ 是单项特征的有理正系数线性组合，则某个倍数 $m\chi$ 本身是单项的；而存在例子（例如 $G=A_5$）使这不可能。

我们概述 Brauer 与 Tate 在 1951 年给出的 Brauer 定理的证明（循 Roquette 的想法）。首先注意：幂零群的任何特征都是单项的。因为若 $G$ 幂零而非交换（交换情形当然平凡），它含一个真包含中心的正规交换子群 $A$。设 $(V,\rho)$ 是 $G$ 的不可约表示，考虑其到 $A$ 的限制。若 $\rho$ 不忠实，可对 $|G|$ 归纳而得结论。若 $\rho$ 忠实，因 $A$ 非中心，$\rho(A)$ 不通过位似作用，故 $A$ 通过若干不同的特征作用在 $V$ 上。$A$ 的每个特征子空间被 $G$ 保持（因 $A$ 在 $G$ 中正规），且 $G$ 传递地置换 $A$-特征子空间（因 $\rho$ 不可约）。但这意味着 $\rho$ 由这些特征子空间之一的稳定子上的某个表示 $\sigma$ 诱导而来，再次归纳即可完成证明。

$G$ 的子群 $H$ 称为 $p$-初等的（$p$ 为素数），若它是 $p$-群与循环群的乘积。这样的群幂零，故其特征全为单项的。只需证明 $G$ 的每个虚拟特征都是初等子群诱导来的特征的 $\mathbb{Z}$-线性组合。更确切地说，考虑映射
\begin{equation}\tag{4.8}\label{c4:eq:4.8}
\operatorname{Ind}_p\colon\bigoplus_{H\colon\ p\text{-初等}}R_H\longrightarrow R_G,
\end{equation}
它是诸群同态 $\operatorname{Ind}^{G}_{H}$（$H$ 取遍 $p$-初等子群）之和。写 $|G|=p^{a}l$，其中 $(p,l)=1$。策略是证明（常函数）$l$ 落在 $\operatorname{Ind}_p$ 的像中。对所有整除 $|G|$ 的 $p$ 施用此结论，便知 $1$ 落在 $\operatorname{Ind}=\sum_p\operatorname{Ind}_p$ 的像中，这正是所要的，因为由投影公式，$\operatorname{Ind}$ 的像是环 $R_G$ 中的理想。

\begin{annote}
此处 $|G|=p^{a}l$ 用的是小写字母 $l$（ell），与全文表示素数的 $\ell$ 只差字体，极易混淆。下文说 $l$ 落在 $\operatorname{Ind}_p$ 的像中，指的是整数 $|G|/p^{a}$，与 $\ell$ 进上同调无关；读到这里请把两种 $l$ 分开。
\par\vspace{0.15em}{\raggedleft 【来源：Serre 前引书 Ch.~10（A）；原书第 82 页该处记号如此，译文照录】\par}
\end{annote}

代替特征的 $\mathbb{Z}$-线性组合，考虑 $A$-线性组合，其中 $A=\mathbb{Z}[\zeta]$，$\zeta$ 为 $|G|$ 次本原单位根。这个环的好处是 $G$ 的所有特征都取值于 $A$，且只需证明 $l$ 落在 $A\otimes\operatorname{Ind}_p$ 的像中。通过显式构造，证明两件事。

(1) $G$ 上每个取值被 $|G|$ 整除的 $\mathbb{Z}$-值类函数都落在 $A\otimes\operatorname{Ind}_p$ 的像中，且此处还可把 $H$ 限制为 $G$ 的循环子群。为证此，先构造常函数 $|G|$，再用投影公式以及我们允许 $A$ 中系数这一事实。

(2) 设 $x$ 是 $G$ 中阶与 $p$ 互素的元素。令 $C=\langle x\rangle$，$P$ 为中心化子 $Z(x)$ 的 $p$-Sylow 子群，$H=C\times P$。则存在 $\psi\in A\otimes R_H$，它是 $\mathbb{Z}$-值的，使 $\psi'=\operatorname{Ind}^{G}_{H}(\psi)$ 满足
\begin{enumerate}[label=(\alph*),leftmargin=2.6em,itemsep=0.2em]
\item $\psi'(x)\not\equiv0\bmod p$；
\item 对阶与 $p$ 互素且不与 $x$ 共轭的元素 $s$，$\psi'(s)=0$。
\end{enumerate}
用 (2) 容易构造 $A\otimes\operatorname{Ind}_p$ 的像中取值于 $\mathbb{Z}$ 的函数 $\chi$，使 $\chi(g)$ 对任何阶与 $p$ 互素的元素 $g$ 不同余于 $0\bmod p$，从而（这很容易）对任何 $g\in G$ 亦然。用 (1) 与 (2) 所得元素的适当线性组合，便得到函数 $l$。

Brauer 定理蕴涵亚纯延拓与函数方程如下。取不可约 $\chi$，把它写成 $\sum n_i\operatorname{Ind}(\chi_i)$，其中诸 $\chi_i$ 是各子群的一维特征，$n_i$ 为整数。则
\begin{equation}\tag{4.9}\label{c4:eq:4.9}
\Lambda(\chi,s)=\prod_{i}\Lambda(\chi_i,s)^{n_i}.
\end{equation}
由第一步，定理对诸 $\chi_i$ 已知，故对 $\chi$ 成立。注意诸 $\chi_i$ 附属于 $k$ 与 $K$ 之间的各个底域 $k_i$。我们无法保证 $n_i$ 为正，这正解释了为何用此方法只能得到亚纯延拓，而非解析延拓。

\begin{annote}
系数 $n_i$ 取整数而非正整数，是 Brauer 方法只给出亚纯延拓的全部原因：$\prod_i\Lambda(\chi_i,s)^{n_i}$ 在 $n_i<0$ 时引入极点，而这些极点未必被别的因子抵消。书中以 $A_5$ 为例说明正系数做不到：若 $\chi$ 是单项特征的有理正系数组合，则某个倍数 $m\chi$ 本身是单项的，而 $A_5$ 没有这样的单项特征。Artin 猜想正是断言这些极点其实都不存在。
\par\vspace{0.15em}{\raggedleft 【来源：Serre 前引书 Ch.~10（A）；Martinet 前引文（A）】\par}
\end{annote}

Brauer 定理的另一个推论是：GRH（对所有 Dedekind zeta 函数的 Riemann 假设）蕴涵每个 $\Lambda(\chi,s)$ 的零点（或极点）都在 $\operatorname{Re}(s)=1/2$ 上，而这无需假设 Artin 猜想。

\section{Artin 猜想}\label{c4:sec:5}

\par\medskip\noindent\textbf{猜想 5.1}\quad 若 $\chi$ 是 $G=\operatorname{Gal}(K/k)$ 的特征且不含平凡特征，则 Artin L 函数 $L(\chi,s)$ 在整个复平面上解析。

如我们所见，此猜想对一维 $\chi$ 成立，更一般地对单项特征成立，因而对幂零群 $G$ 成立；它在函数域情形也成立。此外，既知 $L(\chi,s)$ 亚纯，只需对某个 $m$ 证明关于 $L(m\chi,s)=L(\chi,s)^{m}$ 的猜想。

\par\medskip\noindent\textbf{推论 5.2}\quad 若 $\chi=\sum r_i\operatorname{Ind}^{G}_{H_i}\chi_i$，其中 $0<r_i\in\mathbb{Q}$，诸 $\chi_i$ 一维且非平凡，则 $\chi$ 满足 Artin 猜想。

\par\medskip\noindent\textbf{例 5.1（Aramata--Brauer）}\quad $G$ 的增广表示（正则表示减去平凡表示）就是这种类型。于是对 Galois 的 $K/k$，$\zeta_K(s)/\zeta_k(s)$ 全纯。对非 Galois 的 $K/k$，相应定理仍未解决，尽管它可由猜想推出。

\begin{annote}
推论 5.2 与例 5.1 的差别在于系数 $r_i$ 的正性。Aramata--Brauer 的例子来自增广表示 $\mathrm{Reg}-1$ 的分解：每个不可约表示在正则表示中的重数等于其维数，系数非负，于是 $\zeta_K/\zeta_k$ 全纯。对非 Galois 的 $K/k$，只能写 $\zeta_K/\zeta_k=\prod_{\rho\neq1}L(\rho,s)^{n_\rho}$ 而 $n_\rho\ge0$，但每个 $L(\rho,s)$ 单独的解析性仍未知，商的全纯性因此无从下手。
\par\vspace{0.15em}{\raggedleft 【来源：Martinet 前引文（A）；Serre《Zeta and L functions》（A）】\par}
\end{annote}

Artin 猜想今天已被纳入一般的 Langlands 纲领，并由一般 Langlands 猜想推出。最后一节我们将在映入 $\mathrm{GL}(2)$ 的表示这一情形略作说明。

\section{局部常数}\label{c4:sec:6}

Artin 根数是非常有趣的量。例如，对导子为 $f$ 的 Dirichlet 特征 $\chi$，它由
\begin{equation}\tag{6.1}\label{c4:eq:6.1}
W(\chi)=\frac{\tau(\chi)}{\sqrt{f}\,i^{\delta}}
\end{equation}
给出，其中 $\delta=0$ 或 $1$，视 $\chi(-1)=1$ 或 $-1$ 而定；而 $\tau(\chi)$ 是 Gauss 和
\begin{equation}\tag{6.2}\label{c4:eq:6.2}
\tau(\chi)=\sum_{a=1}^{f}\chi(a)e^{2\pi ia/f}.
\end{equation}

Artin 根数的一些初等性质是：
\begin{itemize}[leftmargin=2.2em,itemsep=0.2em]
\item $W(\chi_1+\chi_2)=W(\chi_1)W(\chi_2)$；
\item $W\bigl(\operatorname{Ind}^{G}_{H}\chi\bigr)=W(\chi)$；
\item $W(\tilde\chi)=\overline{W(\chi)}$。
\end{itemize}

注意最后一条性质的推论：若 $\chi$ 实，则 $\Lambda(\chi,1-s)=W(\chi)\Lambda(\chi,s)$ 且 $W(\chi)=\pm1$。

实特征对应的表示与其反偶同构。这等价于说 $V$ 有非退化的 $G$-不变双线性型。特别地，若存在 $\rho$-不变的实结构 $V_{\mathbb{R}}$，则 $\chi$ 实。众所周知，此时 $\rho$ 是正交的，即可用正交矩阵实现。此时 Fr\"ohlich--Queyrut 的一个定理说 $W(\chi)=1$。

在推导 Hecke L 级数的函数方程时，Tate 发现 $W(\chi)$ 的一个典范分解，即分解为 $k$ 的所有素位上的乘积：
\begin{equation}\tag{6.3}\label{c4:eq:6.3}
W(\chi)=\prod_{v}W_v(\chi).
\end{equation}
$W_v(\chi)$ 只依赖 $\chi$ 到 $G_v$ 的限制（回忆此时 $G$ 交换、$\chi$ 一维）。若 $v$ 阿基米德，$W_v(\chi)=i^{-\delta}$，$\delta=0$ 或 $1$，视 $\chi|_{G_v}$ 平凡与否而定。对非阿基米德 $v$，令
\begin{equation}\tag{6.4}\label{c4:eq:6.4}
\mathfrak{f}_v(\chi)=v^{f(\chi,v)}
\end{equation}
为 $\chi$ 的导子的 $v$ 部分。则
\begin{equation}\tag{6.5}\label{c4:eq:6.5}
W_v(\chi)=\frac{\tau(\chi_v)}{\sqrt{N\mathfrak{f}_v(\chi)}},
\end{equation}
而 $\tau(\chi_v)$ 是联系于特征 $\chi_v\colon k_v^{\times}\to\mathbb{C}^{\times}$ 的某个广义 Gauss 和，后者经由局部类域论对应于 $\chi|_{G_v}$（细节见所引 Tate 的文章）。可与 Dirichlet L 函数的情形相比较。

\begin{annote}
(6.5) 把整体根数拆成局部根数之积，是 Tate 论文的核心副产品：每个 $W_v(\chi)$ 只依赖 $\chi|_{G_v}$，非阿基米德情形由局部 Gauss 和 $\tau(\chi_v)$ 除以 $\sqrt{N\mathfrak{f}_v(\chi)}$ 给出，阿基米德情形退化为 $i^{-\delta}$。非阿基米德情形有 $|\tau(\chi_v)|=(N\mathfrak{f}_v(\chi))^{1/2}$，故每个 $|W_v(\chi)|=1$，整体根数的模也为 $1$；定理 4.1 中“绝对值 $1$”的断言即由此而来。
\par\vspace{0.15em}{\raggedleft 【来源：Tate《Local constants》，Durham 会议录（A）；Deligne《Les constantes des équations fonctionnelles des fonctions L》（A）】\par}
\end{annote}

Langlands（1970）注意到：作为其把 $\mathrm{GL}_n(k)$ 上自守 L 函数与次数 $n$ 的 Artin L 函数相联系的猜想的一个推论，Artin 根数 $W(\chi)$ 在非交换情形也应能表为典范定义的局部根数 $W_v(\chi)$ 之积。他宣布了这一点，未给证明。Deligne（1974）在 Dwork（1956）某些早期结果的基础上给出了一个简短的全局证明，Tate 的叙述即基于 Deligne 的。

把 $\chi$ 表为单项特征的 $\mathbb{Z}$-线性组合的任何方式
\begin{equation}\tag{6.6}\label{c4:eq:6.6}
\chi=\sum_{i}n_i\operatorname{Ind}^{G}_{H_i}\chi_i
\end{equation}
都给出恒等式
\begin{equation}\tag{6.7}\label{c4:eq:6.7}
W(\chi)=\prod_{i}W(\chi_i)^{n_i},
\end{equation}
并提示用 $W_v(\chi_i)$ 类似地定义 $W_v(\chi)$。此处若 $k_i$ 是 $H_i$ 的不动域，则 $W_v(\chi_i)$ 指 Tate 局部常数 $W_w(\chi_i)$ 在 $k_i$ 的所有位于 $v$ 之上的位 $w$ 上的乘积。困难在于证明它与所选表达式无关；对整体根数我们知道此事，只因 $W(\chi)$ 还可由函数方程来定义。事实上这并不完全对，此处正是可延拓性与强可延拓性的区别。然而，若诸 $\chi$ 是维数 $0$ 的虚拟特征（同维数特征的形式差），则它确实成立；而这就够了，因为此时定义
\begin{equation}\tag{6.8}\label{c4:eq:6.8}
W_v(\chi)=W_v(\chi-\dim\chi\cdot1)\,W_v(1)^{\dim\chi}.
\end{equation}
于是证明的主要部分是：若
\begin{equation}\tag{6.9}\label{c4:eq:6.9}
\sum_{i}n_i\operatorname{Ind}^{G}_{H_i}\chi_i=\sum_{j}n'_j\operatorname{Ind}^{G}_{H'_j}\chi'_j,
\end{equation}
则
\begin{equation}\tag{6.10}\label{c4:eq:6.10}
\prod_{i}W_v(\chi_i)^{n_i}W_v(1_i)^{-n_i}
=\prod_{j}W_v(\chi'_j)^{n'_j}W_v(1_j)^{-n'_j}.
\end{equation}
此方程中出现的局部根数是交换的，而 $1_i$ 是 $k_i$（即 $H_i$ 的不动域）的平凡特征。

{\small
\par\medskip\noindent\textbf{参考文献}\par\nopagebreak
\noindent[1] J. Tate, Local constants, in \emph{Algebraic Number Fields: Proceedings of the Durham Symposium}, Academic Press, New York, 1977.\par
}

\section{\texorpdfstring{与模形式的关系（$\mathrm{GL}(2)$ 情形）}{与模形式的关系（GL(2) 情形）}}\label{c4:sec:7}

Artin 猜想是检验 motivic L 函数与自守 L 函数之间关系的最早场合之一。尽管至今结果仍是部分的，我们不能在本综述中略去这一主题。我们假定读者熟悉 $\mathbb{Q}$ 上 $\mathrm{GL}(2)$ 的经典模形式。

本节中 $k=\mathbb{Q}$，我们固定表示 $\rho\colon\operatorname{Gal}(K/\mathbb{Q})\to GL(2,\mathbb{C})$，并自始至终假定它是奇的，即 $\epsilon=\det(\rho)$ 把复共轭送到 $-1$。注意 $\epsilon$ 可看作奇的 Dirichlet 特征。

模形式是 Galois 表示的一个非常重要的来源。历史上，附着于权 $2$ 本征形式的 $\ell$ 进表示先被发现（Eichler 与 Shimura），然后是附着于权 $\ge2$ 本征形式的表示（Deligne）。直到 1974 年，Deligne 与 Serre 才找到附着于权 $1$ 本征形式的较简单的复表示。

\par\medskip\noindent\textbf{定理 7.1（Deligne--Serre）}\quad 设
\begin{equation}\tag{7.1}\label{c4:eq:7.1}
f=\sum_{n=1}^{\infty}a_nq^{n}\qquad(q=e^{2\pi i\tau})
\end{equation}
是 $\Gamma_0(N)$ 上的正规化（$a_1=1$）权 $1$ 新形式，nebentypus 为 $\epsilon$（模 $N$ 的奇 Dirichlet 特征）。则存在不可约表示 $\rho\colon\operatorname{Gal}(K/\mathbb{Q})\to GL(2,\mathbb{C})$，其导子为 $N$、行列式为 $\epsilon$，且其 Artin L 函数是
\begin{equation}\tag{7.2}\label{c4:eq:7.2}
L(\rho,s)=L(f,s)=\sum a_nn^{-s}.
\end{equation}

这些表示满足 Artin 猜想，因为已知 $L(f,s)$ 在整个复平面上解析。

Deligne--Serre 的证明大致如下。为简单起见，设 $f$ 的 Fourier 系数在 $\mathbb{Z}$ 中。我们的目的是构造表示，使 $\operatorname{tr}(\rho(\sigma_p))=a_p$ 且 $\det(\rho(\sigma_p))=\epsilon(p)$。权 $2$ 本征形式附着的表示已相当清楚，它们来自模曲线的 Jacobian 簇的 $\ell$ 进 Tate 模。把它们模 $\ell$ 约化，可得 $\mathrm{GL}(2,\mathbb{F}_\ell)$ 中的表示。Deligne 与 Serre 从 $f$ 出发，把它乘以某个权 $1$ 的 Eisenstein 级数（其 $q$-展开模 $\ell$ 同余于 $1$），得到权 $2$ 本征形式，其 $q$-展开模 $\ell$ 同余于 $f$ 的 $q$-展开。于是他们构造出的还不是所要的表示 $\rho$，而是一个“模 $\ell$”的近似：表示 $\bar\rho_\ell\colon\operatorname{Gal}(\overline{\mathbb{Q}}/\mathbb{Q})\to\mathrm{GL}(2,\mathbb{F}_\ell)$，满足 $\operatorname{tr}(\bar\rho_\ell(\sigma_p))\equiv a_p\bmod\ell$，$\det(\bar\rho_\ell(\sigma_p))\equiv\epsilon(p)\bmod\ell$。他们对无穷多个 $\ell$ 做到了这一点。

下一步是用 Rankin 的一个解析结果证明 $a_p$ 只有有限多个，至少当我们排除一个密度（任意）小的素数集时如此。权 $1$ 的关键性正在此处。用此结果得到 $\bar\rho_\ell$ 的像的一致界。

最后一步是把诸 $\bar\rho_\ell$“粘合”起来，得到映入 $\mathrm{GL}(2,\mathcal{O}_E)$（$\mathcal{O}_E$ 为某个整数环）的表示 $\rho$，它在无穷多个 $\ell$ 处约化为 $\bar\rho_\ell$。由于第二步所得的 $\bar\rho_\ell$ 的像的界，这是可能的。此后证明 $\rho$ 具有所要的性质就相对容易了。

反方向上，有下述重要定理，其（在此情形）归于 A.~Weil，可追溯到 1967 年。它称为“Weil 逆定理”，是同类型更多定理的第一例。

\par\medskip\noindent\textbf{定理 7.2（Weil）}\quad 设 $\rho$ 是 $\operatorname{Gal}(\overline{\mathbb{Q}}/\mathbb{Q})$ 的奇、不可约、二维复表示。设 $N$ 为其导子，$\epsilon$ 为其行列式。假设存在 $M$，使对每个导子与 $M$ 互素的 Dirichlet 特征 $\chi$，Artin L 函数 $L(\rho\otimes\chi,s)$ 在整个复平面上解析。则存在级 $N$、权 $1$、nebentypus 为 $\epsilon$ 的正规化新形式，使
\begin{equation}\tag{7.3}\label{c4:eq:7.3}
L(\rho,s)=L(f,s).
\end{equation}

此定理证明的一个极好来源是 Ogg 的书。

这两个定理指向 Artin 猜想与 Galois 表示模性问题之间的紧密联系。事实上，它们表明：对本节的 $\rho$，两个问题等价。有鉴于此，一些作者把“$\rho$ 来自上述 $f$”这一断言称为强 Artin 猜想。

下一个大步由 Langlands 迈出，并由 Tunnell 补充\footnote{原书作“Tunnel”，标准拼法为 Tunnell。}。$\mathrm{GL}(2,\mathbb{C})$ 的有限子群 $G$ 按其像 $\tilde G$ 在 $\mathrm{PGL}(2,\mathbb{C})$ 中的情形分类。Jordan 的一个经典定理说 $\tilde G$ 可以是循环群、二面体群、$A_4$（四面体情形）、$S_4$（八面体情形）或 $A_5$（二十面体情形）。

下一个定理非常深刻，即使一个梗概（见所引 Gelbart 的论文）也超出本综述的范围。

\par\medskip\noindent\textbf{定理 7.3（Langlands--Tunnell）}\quad 设 $\rho$ 奇、不可约、二维，且 $G=\operatorname{Im}(\rho)$ 属于上面列出的前四种类型之一（$\tilde G$ 循环、二面体、四面体或八面体——事实上循环情形不会出现，因为此时 $\rho$ 不能不可约）。则 $L(\rho,s)$ 满足强 Artin 猜想，即 $\rho$ 是模的。

我们以如下说明结束这篇极为简略且不完整的 Artin L 函数综述：二维表示的 Artin 猜想的二十面体情形近年来由 R.~Taylor 及其合作者研究。关于 Taylor 纲领的简短叙述（写于大部分结果获得之前）见 R.~Taylor 的 \emph{Icosahedral Galois Representations}，\emph{Pacific J. Math.}，1997，337--347。

\begin{annote}
定理 7.3 覆盖了除二十面体以外的全部类型。二十面体情形之所以难，在于 $A_5$ 非可解，Langlands--Tunnell 所用的那套提升依赖可解性——先限制到可解子群，再基变换回原群——而这一套在 $A_5$ 上失效。Taylor 的纲领改用 $\ell$ 进上同调与自守提升处理它，至今仍只有部分结果，这正是本章末尾所引那篇 1997 年短文写作时的状况。
\par\vspace{0.15em}{\raggedleft 【来源：Gelbart《Three lectures on the modularity of $\bar\rho_{E,3}$ and the Langlands reciprocity conjecture》（A）；Taylor《Icosahedral Galois Representations》，Pacific J.~Math.（A）】\par}
\end{annote}

{\small
\par\medskip\noindent\textbf{参考文献}\par\nopagebreak
\noindent[1] J.-P. Serre, Modular forms and Galois representations, \emph{Algebraic Number Fields: Proceedings of the Durham Symposium}, Academic Press, New York, 1977.\par
\noindent[2] A. Ogg, \emph{Modular Forms and Dirichlet Series}, Benjamin, New York, 1969.\par
\noindent[3] S. Gelbart, Three lectures on the modularity of $\bar\rho_{E,3}$ and the Langlands reciprocity conjecture, in G.~Cornell, J.~H. Silverman, and G.~Stevens, eds., \emph{Modular Forms and Fermat's Last Theorem}, Springer-Verlag, New York, 1997.\par
}

\let\thesection\chFourOldTheSection
